\chapter{Исследовательский раздел}\label{research}

\section{Цель исследования}

Целью исследования является оценка эффективности разработанного метода синхронизации состояний двух игроков и его сравнение с базовой событийной синхронизацией без прикладной надежности.

В рамках исследования необходимо:
\begin{itemize}
	\item определить критерии оценки эффективности синхронизации состояний;
	\item задать параметры сетевой среды для экспериментальных запусков;
	\item выполнить серию испытаний для разработанного метода и базового UDP-подхода;
	\item оценить долю успешно примененных событий;
	\item оценить среднее время подтверждения или возврата события;
	\item оценить количество повторных отправок, возникающих при потере сообщений;
	\item определить влияние дублирования пакетов на повторную обработку событий;
	\item сравнить устойчивость двух подходов к потере, задержке и дублированию UDP-пакетов.
\end{itemize}

\section{Сравниваемые режимы синхронизации}

Для проведения исследования в разработанном программном обеспечении реализованы два режима синхронизации, выбираемые параметром \texttt{SynchronizationResearchMode}. Режим \texttt{ReliableFsmMethod} соответствует разработанному методу. Режим \texttt{BasicUdpEvents} используется как контрольный вариант и моделирует простую событийную передачу по UDP без подтверждений и повторной отправки.

Сравнение выполняется при одинаковых входных событиях и одинаковых параметрах сетевой среды. Это позволяет исключить влияние различий в игровой логике и сравнивать именно сетевые и алгоритмические механизмы синхронизации.

\begin{table}[H]
	\caption{Сравниваемые режимы синхронизации}
	\label{tab:research-modes}
	\centering
	\small
	\renewcommand{\arraystretch}{1.2}
	\begin{tabular}{|p{3.4cm}|p{4.6cm}|p{4.6cm}|}
		\hline
		Механизм & Базовая событийная UDP-синхронизация & Разработанный метод \\ \hline
		Режим синхронизации & \texttt{BasicUdpEvents} & \texttt{ReliableFsmMethod} \\ \hline
		Передача события & Однократная отправка события & Отправка события с регистрацией ожидания подтверждения \\ \hline
		Подтверждение обработки & Не используется & Используется сообщение \texttt{Acknowledgement} \\ \hline
		Повторная отправка & Не выполняется & Выполняется после истечения \texttt{resendDelay} \\ \hline
		Фильтрация дубликатов & Не выполняется на уровне метода & Выполняется по ключу события \\ \hline
		Контроль порядка & Не выполняется & Выполняется по \texttt{sequenceNumber} \\ \hline
		Защита состояния FSM & Не выполняется & Обеспечивается проверкой порядка, подтверждением и повторной отправкой \\ \hline
	\end{tabular}
\end{table}

В базовом режиме событие, потерянное на транспортном уровне, не восстанавливается. Если сообщение было продублировано, оно может быть повторно передано в обработчик конечного автомата. В разработанном методе потеря события или подтверждения приводит к повторной отправке, а повторное получение уже обработанного события не изменяет состояние конечного автомата.

\section{Критерии оценки эффективности}

Для оценки используются метрики, собираемые при выполнении разработанного программного обеспечения классом \texttt{ResearchMetrics}. Метрики фиксируются отдельно для каждого запуска и выводятся в журнал приложения после завершения автоматического теста.

К основным метрикам относятся:
\begin{itemize}
	\item $N_{created}$ --- количество созданных логически значимых событий;
	\item $N_{sent}$ --- количество фактически отправленных сетевых сообщений;
	\item $N_{lost}$ --- количество имитированных потерь исходящих UDP-пакетов;
	\item $N_{resend}$ --- количество повторных отправок событий;
	\item $N_{dup}$ --- количество отклоненных дубликатов;
	\item $N_{wrong}$ --- количество событий, отклоненных из-за нарушения порядка;
	\item $N_{applied}$ --- количество событий, примененных к конечному автомату;
	\item $T_{ack}^{avg}$ --- среднее время подтверждения или возврата события.
\end{itemize}

Для разработанного метода величина $T_{ack}^{avg}$ соответствует среднему времени получения подтверждения \texttt{Acknowledgement}. Для базового режима эта величина интерпретируется как среднее время возврата события от серверной стороны к клиенту, поскольку отдельное сообщение подтверждения в данном режиме не используется.

Среднее время подтверждения или возврата события определяется выражением:
\[
	T_{ack}^{avg} = \frac{1}{N}\sum_{i=1}^{N}(t_i^{ack} - t_i^{send}),
\]
где:
\begin{itemize}
	\item $N$ --- количество событий, для которых был получен ответ;
	\item $t_i^{send}$ --- момент создания события;
	\item $t_i^{ack}$ --- момент получения подтверждения или подтвержденного сервером события.
\end{itemize}

Доля примененных событий определяется как:
\[
	K_{applied} = \frac{N_{applied}}{N_{created}}.
\]

Для разработанного метода ожидаемым значением $K_{applied}$ является значение, близкое к единице. Для базового режима при потере пакетов значение $K_{applied}$ может уменьшаться, так как потерянные события не восстанавливаются. При дублировании пакетов в базовом режиме показатель может превышать единицу, что указывает на повторную обработку одного и того же события.

Относительное количество повторных отправок определяется как:
\[
	K_{resend} = \frac{N_{resend}}{N_{created}}.
\]

Рост $K_{resend}$ при наличии потерь является ожидаемым для разработанного метода, так как дополнительные сообщения используются для восстановления доставки критичных событий. 
Для базового UDP-подхода $K_{resend}$ равно нулю, однако отсутствие повторов в этом случае не является преимуществом, так как потерянные события не компенсируются.

\section{Порядок проведения исследования}

Автоматизация испытаний выполняется компонентом \texttt{AutomatedResearchTest}. Компонент создает 100 событий с интервалом 0.2~с. После завершения генерации событий выполняется ожидание завершения сетевой обработки, затем компонент \texttt{ResearchMetrics} выводит итоговые значения метрик.

Параметры сетевой среды задаются в компоненте \texttt{UdpEventTransport}:
\begin{itemize}
	\item \texttt{outgoingPacketLossChance} --- вероятность потери исходящего пакета;
	\item \texttt{artificialLatencyMs} --- искусственная задержка отправки;
	\item \texttt{jitterMs} --- вариативность задержки;
	\item \texttt{duplicatePacketChance} --- вероятность дублирования отправляемого пакета.
\end{itemize}

Для каждого сценария выполняются два запуска: один в режиме \texttt{ReliableFsmMethod}, второй в режиме \texttt{BasicUdpEvents}. Значения параметров сетевой среды в двух запусках совпадают.

\section{Сценарии исследования}

Для исследования выбраны сценарии, отражающие типовые условия сетевого взаимодействия в многопользовательском приложении. Сценарии приведены в таблице~\ref{tab:research-scenarios}.

\begin{table}[H]
	\caption{Сценарии исследования UDP-обмена}
	\label{tab:research-scenarios}
	\centering
	\renewcommand{\arraystretch}{1.2}
	\begin{tabular}{|p{4cm}|c|c|c|c|}
		\hline
		Сценарий & Потери, \% & Задержка, мс & Jitter, мс & Дубликаты, \% \\ \hline
		Стабильная сеть & 0 & 0 & 0 & 0 \\ \hline
		Умеренная задержка & 0 & 100 & 20 & 0 \\ \hline
		Потери пакетов & 5 & 100 & 20 & 0 \\ \hline
		Высокие потери & 15 & 100 & 30 & 0 \\ \hline
		Дублирование пакетов & 0 & 100 & 20 & 20 \\ \hline
		Нестабильная сеть & 10 & 150 & 50 & 10 \\ \hline
	\end{tabular}
\end{table}

Первые два сценария используются для оценки влияния задержки без потери сообщений. Сценарии с потерями позволяют оценить способность метода восстанавливать доставку критичных событий. Сценарий с дублированием проверяет устойчивость к повторному получению одного и того же события. Нестабильная сеть объединяет задержку, вариативность задержки, потерю и дублирование пакетов.

\section{Результаты исследования разработанного метода}

В таблице~\ref{tab:research-reliable-results} приведены результаты запусков для режима \texttt{ReliableFsmMethod}. 

\begin{table}[H]
	\caption{Результаты исследования разработанного метода}
	\label{tab:research-reliable-results}
	\centering
	\small
	\renewcommand{\arraystretch}{1.2}
	\begin{tabular}{|p{3.2cm}|c|c|c|c|c|c|}
		\hline
		Сценарий & Создано & Отправлено & Потери & Повторы & Применено & $T_{ack}^{avg}$, с \\ \hline
		Стабильная сеть & 100 & 102 & 0 & 2 & 100 & 0.099 \\ \hline
		Умеренная задержка & 100 & 118 & 0 & 18 & 100 & 0.332 \\ \hline
		Потери пакетов & 100 & 146 & 8 & 54 & 100 & 0.939 \\ \hline
		Высокие потери & 100 & 132 & 20 & 46 & 99 & 0.711 \\ \hline
		Дублирование пакетов & 100 & 131 & 0 & 12 & 100 & 0.362 \\ \hline
		Нестабильная сеть & 100 & 139 & 8 & 38 & 99 & 0.548 \\ \hline
	\end{tabular}
\end{table}

Из таблицы видно, что при стабильной сети метод обеспечивает применение всех созданных событий. При увеличении задержки возрастает среднее время подтверждения, 
однако доля примененных событий остается равной единице. При появлении потерь возрастает количество повторных отправок, что является следствием работы механизма восстановления доставки. 
В сценариях с высокими потерями и нестабильной сетью часть событий может не попасть в окно наблюдения, однако механизм повторной отправки продолжает работу до получения подтверждения.

Важным результатом является то, что увеличение числа отправленных сообщений не приводит к повторному изменению состояния конечного автомата. 
Повторные отправки и дубликаты распознаются по ключу события и используются только для восстановления факта доставки и подтверждения обработки.

\section{Результаты исследования базового UDP-подхода}

В таблице~\ref{tab:research-basic-results} приведены результаты запусков для режима \texttt{BasicUdpEvents}. В этом режиме событие отправляется один раз, не помещается в журнал неподтвержденных событий и не отправляется повторно.

\begin{table}[H]
	\caption{Результаты исследования базовой событийной UDP-синхронизации}
	\label{tab:research-basic-results}
	\centering
	\small
	\renewcommand{\arraystretch}{1.2}
	\begin{tabular}{|p{3.2cm}|c|c|c|c|c|c|}
		\hline
		Сценарий & Создано & Отправлено & Потери & Повторы & Применено & $T_{ack}^{avg}$, с \\ \hline
		Стабильная сеть & 100 & 100 & 0 & 0 & 100 & 0.086 \\ \hline
		Умеренная задержка & 100 & 100 & 0 & 0 & 100 & 0.301 \\ \hline
		Потери пакетов & 100 & 95 & 5 & 0 & 91 & 0.314 \\ \hline
		Высокие потери & 100 & 85 & 15 & 0 & 73 & 0.327 \\ \hline
		Дублирование пакетов & 100 & 120 & 0 & 0 & 124 & 0.309 \\ \hline
		Нестабильная сеть & 100 & 101 & 10 & 0 & 88 & 0.423 \\ \hline
	\end{tabular}
\end{table}

При отсутствии потерь базовый подход способен доставить все события, так как сетевой канал не создает отказов. Однако при потере пакетов доля примененных событий снижается, поскольку потерянное событие не восстанавливается. При дублировании пакетов количество применений может превышать число созданных событий, что свидетельствует о повторной обработке одного и того же логического действия.

\section{Сравнение результатов}

Сводное сравнение двух режимов приведено в таблице~\ref{tab:research-comparison}. 
Для сравнения используется доля примененных событий $K_{applied}$. Для разработанного метода значение, близкое к единице, означает доставку и обработку созданных событий. 
Для базового подхода значение меньше единицы означает потерю событий, а значение больше единицы --- повторную обработку из-за дублирования.

\begin{table}[H]
	\caption{Сравнение разработанного метода с базовой UDP-синхронизацией}
	\label{tab:research-comparison}
	\centering
	\small
	\renewcommand{\arraystretch}{1.2}
	\begin{tabular}{|p{5cm}|c|c|}
		\hline
		Сценарий & $K_{applied}$ базового подхода & $K_{applied}$ разработанного метода \\ \hline
		Стабильная сеть & 1.00 & 1.00 \\ \hline
		Умеренная задержка & 1.00 & 1.00 \\ \hline
		Потери пакетов & 0.91 & 1.00 \\ \hline
		Высокие потери & 0.73 & 0.99 \\ \hline
		Дублирование пакетов & 1.24 & 1.00 \\ \hline
		Нестабильная сеть & 0.88 & 0.99 \\ \hline
	\end{tabular}
\end{table}

Результаты показывают, что при благоприятных сетевых условиях оба подхода дают сопоставимый результат. Это ожидаемо, поскольку при отсутствии потерь 
и дубликатов простая передача события достаточна для сохранения состояния. Отличие проявляется при появлении сетевых ошибок.

При потере пакетов базовый подход теряет часть событий без возможности восстановления. 
Разработанный метод компенсирует такие потери повторной отправкой. При дублировании пакетов базовый подход допускает повторную обработку события, 
тогда как разработанный метод распознает уже обработанное событие и не изменяет конечный автомат повторно.

\section{Выводы по результатам исследования}

По результатам исследования можно сделать следующие выводы:
\begin{itemize}
	\item при стабильной сети разработанный метод и базовая событийная UDP-синхронизация обеспечивают применение всех событий;
	\item при увеличении задержки возрастает среднее время подтверждения, однако согласованность состояний сохраняется;
	\item при потере пакетов базовый подход теряет часть событий, так как не содержит механизма восстановления доставки;
	\item разработанный метод компенсирует потерю сообщений за счет хранения событий в \texttt{pendingEvents} и повторной отправки после истечения времени ожидания;
	\item при дублировании пакетов базовый подход может привести к повторной обработке одного и того же события;
	\item фильтрация дубликатов и контроль порядка в разработанном методе предотвращают повторное и несогласованное изменение конечного автомата;
	\item увеличение числа отправленных сообщений в разработанном методе является ожидаемым следствием обеспечения прикладной надежности поверх UDP.
\end{itemize}

\section*{Вывод}

В данном разделе были определены критерии оценки эффективности разработанного метода, описаны режимы экспериментального сравнения и заданы сценарии исследования UDP-сетевого обмена. Было проведено сравнение разработанного метода с базовой событийной UDP-синхронизацией без прикладной надежности.

Результаты исследования показывают, что базовый подход достаточен только при отсутствии сетевых ошибок. При потере и дублировании пакетов он не гарантирует ни доставку критичного события, ни его однократную обработку. Разработанный метод вводит дополнительные служебные механизмы: подтверждения, повторную отправку, контроль порядка и фильтрацию дубликатов. Эти механизмы увеличивают количество сетевых сообщений, однако позволяют сохранять согласованность состояний конечного автомата в условиях ненадежного UDP-обмена.
